% Ad Atticum 9.15 (ad-atticum-09-15)
% Source: https://raw.githubusercontent.com/PerseusDL/canonical-latinLit/master/data/phi0474/phi057/phi0474.phi057.perseus-lat1.xml
% Fetched by scripts/fetch_latin.py

% Dateline (Perseus): Scr. in Formiano viii K. Apr. a. 70; (49).

\ciceroLetterOpener{CICERO ATTICO salutem}

\ciceroSection{1} cum dedissem ad te litteras ut scires Caesarem Capuae vii Kal. fore, adlatae mihi Capua sunt †et hoc mihi et† in Albano apud Curionem v K. fore. eum cum videro, Arpinum pergam. si mihi veniam quam peto dederit, utar illius condicione; si minus, impetrabo aliquid a me ipso. ille , ut ad me scripsit, legiones singulas posuit Brundisi, Tarenti, Siponti. claudere mihi videtur maritimos exitus et tamen ipse Graeciam spectare potius quam Hispanias.

\ciceroSection{2} sed haec longius absunt. me nunc et congressus huius stimulat (is vero adest) et primas eius actiones horreo. volet enim, credo, S. C. facere, volet augurum decretum (rapiemur aut absentes vexabimur), vel ut consules roget praetor vel dictatorem dicat; quorum neutrum ius est. etsi , si Sulla potuit efficere ab interrege ut dictator diceretur et magister equitum , cur hic non possit? nihil expedio nisi ut aut ab hoc tamquam Q. Mucius aut ab illo tamquam L. Scipio.

\ciceroSection{3} cum tu haec leges, ego illum fortasse convenero. τέτλαθι κύντερον ne illud quidem nostrum proprium. erat enim spes propinqui reditus, erat hominum querela. nunc exire cupimus, qua spe reditus mihi quidem numquam in mentem venit. non modo autem nulla querela est municipalium hominum ac rusticorum sed contra metuunt ut crudelem, iratum. nec tamen mihi quicquam est miserius quam remansisse nec optatius quam evolare non tam ad belli quam ad fugae societatem. † sed tu omnia qui† consilia differebas in id tempus cum sciremus quae Brundisi acta essent. scimus nempe; haeremus nihilo minus. vix enim spero mihi hunc veniam daturum, etsi multa adfero iusta ad impetrandum. sed tibi omnem illius meumque sermonem omnibus verbis expressum statim mittam.

\ciceroSection{4} tu nunc omni amore enitere ut nos cura tua et prudentia iuves. ita subito accurrit ut ne T. Rebilum quidem, ut constitueram, possim videre; omnia nobis imparatis agenda. sed tamen ἄλλα μὲν αὐτόσ , ut ait ille, ἄλλα δὲ καὶ δαίμων ὑποθήσεται . quicquid egero continuo scies. mandata Caesaris ad consules et ad Pompeium quae rogas, nulla habeo †et descripta attulit illa e via†misi ad te ante; e quibus mandata intellegi posse. Philippus Neapoli est, Lentulus Puteolis. de Domitio, ut facis, sciscitare ubi sit, quid cogitet.

\ciceroSection{5} quod scribis asperius me quam mei patiantur mores de Dionysio scripsisse, vide quam sim antiquorum hominum. te medius fidius hanc rem gravius putavi laturum esse quam me. nam praeter quam quod te moveri arbitror oportere iniuria quae mihi a quoquam facta sit, praeterea te ipsum quodam modo hic violavit cum in me tam improbus fuit. sed tu id quanti aestimes tuum iudicium est; nec tamen in hoc tibi quicquam oneris impono. ego autem illum male sanum semper putavi, nunc etiam impurum et sceleratum puto nec tamen mihi inimiciorem quam sibi. Philargyro bene curasti. causam certe habuisti et veram et bonam, relictum esse me potius quam reliquisse.

\ciceroSection{6} cum dedissem iam litteras a. d. viii Kal., pueri quos cum Matio et Trebatio miseram epistulam mihi attulerunt hoc exemplo: MATIVS ET TREBATIVS CICERONI IMP. salutem cum Capua exissemus, in itinere audivimus Pompeium Brundisio a. d. xvi K. Aprilis cum omnibus copiis quas habuerit profectum esse; Caesarem postero die in oppidum introisse, contionatum esse, inde Romam contendisse, velle ante K. esse ad urbem et pauculos dies ibi commorari, deinde in Hispanias proficisci. nobis non alienum visum est, quoniam de adventu Caesaris pro certo habebamus, pueros tuos ad te remittere, ut id tu quam primum scires. mandata tua nobis curae sunt eaque ut tempus postularit agemus. Trebatius sedulo facit ut antecedat. epistula conscripta nuntiatum est nobis Caesarem a. d. viii K. Aprilis Beneventi mansurum, a. d. vii Capuae, a. d. vi Sinuessae. hoc pro certo putamus.
